\chapter{相关技术}\label{chap:tech}

本章介绍 QuickBoard 系统实现所涉及的关键技术与选型依据，涵盖前端框架与构建工具、画布渲染、实时通信、公式渲染与 OCR 服务等。

\section{前端框架与工程化}
\subsection{Vue3 与 Vite}
QuickBoard 前端采用 Vue3 进行组件化开发，通过响应式状态管理实现工具状态、房间信息与 UI 交互的联动。构建工具选择 Vite，以更快的开发热更新与生产构建能力提升迭代效率。该组合能够在保证工程可维护性的同时降低构建与调试成本。\cite{vue3,vite}

\subsection{Tailwind CSS}
项目使用 Tailwind CSS 进行样式组织，通过原子类组合实现按钮、输入框、卡片等基础样式，并结合少量自定义样式定义动效与背景纹理，从而在不引入复杂设计系统的情况下满足极简风格需求。

\section{画布渲染与交互：Fabric.js}
在线白板的核心是“可编辑画布”。QuickBoard 使用 Fabric.js 封装原生 Canvas 提供的绘制能力，将图形元素抽象为对象（Object），并提供选择、拖拽、缩放、命中检测等交互能力。相比直接操作 Canvas API，Fabric 更易于管理多对象场景，并支持序列化与反序列化，便于协作同步与导出。

在性能方面，项目通过视口裁剪与对象缓存等策略降低不必要的重绘开销，同时在选中对象时关闭缓存以避免编辑过程中的视觉滞后。\cite{fabricjs}

\subsection{Canvas 性能特征与优化方向}
Canvas 渲染通常运行在单线程主线程上，渲染与事件处理共享同一执行栈。当对象数量增大时，性能瓶颈往往来自三个方面：对象遍历成本（每帧遍历大量对象）、不可见对象绘制成本（视口外对象仍参与绘制）、以及重复绘制成本（静态对象每帧被重复光栅化）。因此，优化方向也相对明确：通过视口裁剪跳过不可见对象；通过对象缓存对静态对象复用位图；通过减少无效重绘与合并绘制批次降低每帧开销。QuickBoard 的渲染优化围绕上述方向展开，并以基准测试量化优化前后的平均渲染耗时变化。

\section{实时通信：Socket.IO}
实时协作需要低延迟、双向通信与房间管理机制。QuickBoard 后端基于 Socket.IO 提供房间语义、断线重连与事件广播能力；前端通过 Socket.IO Client 进行连接、加入房间、发送绘制事件与接收他人更新。相较于直接使用原生 WebSocket，Socket.IO 在工程实践中提供了更完善的事件模型与连接管理能力，降低了实现复杂度。\cite{socketio}

\subsection{WebSocket 协议基础}
浏览器与服务端的实时通信通常建立在 WebSocket 协议之上。WebSocket 通过一次 HTTP 握手升级为全双工长连接，减少轮询带来的额外开销，并能以较低延迟推送事件\cite{rfc6455}。在白板协作场景中，绘制增量、对齐事件与成员状态变更都具有“频繁、小包、时效敏感”的特点，长连接模型能够更好地匹配这种事件流。

\section{同步策略与冲突处理}
多人编辑不可避免地带来并发冲突与消息乱序问题。QuickBoard 采用“服务端权威裁决 + 增量同步”的设计思想：客户端发送操作增量，服务端对关键字段进行最后写入优先（LWW）裁决后再广播被接受的变更；同时维护房间版本号与增量日志，使新加入或断线重连的客户端可以通过版本差量补齐状态，必要时再回退到全量快照。该策略在工程复杂度与一致性体验之间取得平衡，适合毕业设计的轻量协作场景。

\subsection{OT 与 CRDT 的对比与取舍}
协作一致性领域常见的两条技术路线为 OT 与 CRDT。OT 通过对并发操作进行转换来保持意图，并在早期协同编辑系统中得到广泛应用\cite{sun_ellis_1998}；CRDT 则以“可交换、可合并的数据结构”为核心，天然支持最终一致性与离线合并\cite{shapiro_rr7506_2011}。二者的共同点是能够在并发修改下保证收敛，但在工程实现成本、调试复杂度与数据模型适配性方面存在差异。

QuickBoard 的白板状态以“图形对象集合”为核心：对象数量多、字段多、交互频繁，且系统定位为短会话协作。若引入通用 CRDT 框架，需要为对象集合与字段更新设计更复杂的可交换结构，并处理 tombstone、压缩与持久化等配套问题；若采用 OT，需要为各类对象操作定义转换规则并在服务端维护操作序列。综合考虑毕业设计的交付周期与可解释性要求，QuickBoard 采用更轻量的字段级 LWW，并以服务端权威裁决降低分叉概率，通过版本补包与快照兜底保证最终收敛。

\subsection{逻辑时钟与 LWW 合并}
在 LWW 策略中，“写入的新旧”需要一个可比较的时间戳。若直接使用客户端物理时间，可能因时钟偏差导致误判；因此工程上常使用逻辑时钟为更新分配单调递增的序号，并在接收远端更新后对本地逻辑时钟进行追赶。该思路与分布式系统中对事件顺序的经典讨论一致\cite{lamport_1978}。QuickBoard 将时间戳关联到对象字段，以字段级别进行比较与合并，从而避免“整对象覆盖”带来的误更新与额外网络开销。

\section{公式渲染与编辑：KaTeX}
为了支持数学公式的高质量排版，项目在前端使用 KaTeX 将 LaTeX 字符串渲染为数学公式。与将 LaTeX 作为纯文本绘制到 Canvas 不同，QuickBoard 采用渲染层叠加的方式显示公式，并提供双击编辑、实时预览与确认插入流程，从而提升可读性与可编辑性。\cite{katex}

\subsection{可编辑渲染与对象模型的协调}
在白板对象模型中，公式既需要像普通对象一样被选择、移动、缩放与序列化，又需要在视觉上以高质量排版呈现。QuickBoard 将公式的“数据层”和“显示层”分离：数据层以 LaTeX 字符串形式作为对象字段参与同步与持久化；显示层由 KaTeX 在前端渲染生成可读的排版结果。该分离使得跨端同步时传输的是稳定的文本表示，而不是与分辨率或字体强相关的位图；同时也为“二次编辑”提供了自然入口——用户编辑的是 LaTeX 数据，渲染层随之更新。

\section{数学公式识别：OCR 服务}
QuickBoard 的公式识别功能采用独立 OCR 服务实现。前端对白板选区进行截图并发送到后端接口，后端将请求转发至 OCR 服务并返回识别结果。OCR 服务侧对输入图像进行必要的预处理（如尺寸对齐、内容裁剪、反色/二值化等），并在结果为空或包含噪声时进行清洗与回退尝试，最终以统一格式返回 LaTeX 或错误信息。该“前端采集—后端代理—OCR 服务识别”的三段式架构便于部署与调试，也降低了前端复杂度。\cite{fastapi}

\subsection{预处理的重要性}
数学 OCR 的输入具有明显的“非规范性”：笔迹粗细与倾斜程度差异大，背景可能包含网格或色块，框选区域可能过紧导致字符被裁剪，也可能过松引入大量噪声。若直接将原始截图输入识别模型，空结果或错误结果的概率显著上升。因此，工程实践中通常需要在服务端加入预处理步骤，例如对输入尺寸进行对齐、对深色背景做反色、对前景进行二值化或增强，并在必要时构造多种预处理候选以提升命中概率。QuickBoard 将预处理与失败回退作为 OCR 链路的核心组成部分，并通过统一错误码与 traceId 让失败原因具备可定位性。

\subsection{错误码与可定位性}
与“返回空字符串”相比，统一错误码能够把失败原因分类为可操作的工程问题：例如超时（资源或网络问题）、空结果（输入质量或模型能力问题）、坏图（裁剪或编码问题）。traceId 则提供跨组件追踪的锚点，使一次识别请求可以在前端日志、后端代理与 OCR 服务端日志中被关联。通过错误码与 traceId，OCR 评估可以从“好不好用”的主观感受，落到“哪些失败类型占主导、如何改进”的工程分析路径上。

\section{极简白板产品参考}
为提升“快速进入、低干扰”的首屏体验，首页信息架构参考了 Excalidraw 等极简白板产品的设计取向，并结合房间化协作入口进行本地化改造。\cite{excalidraw}
